\begin{textsample}{NEG Dim 2 – 2018\_c5 – Score -2.00 – 2018\_c5/enade\_2018\_admin\_6a5e0b186f.txt}  \label{ex:f2_neg_114}
CAPÍTULO 5 DISTRIBUIÇÃO DOS CONCEITOS 5.1 PANORAMA NACIONAL DA DISTRIBUIÇÃO DOS CONCEITOS O Conceito Enade é um indicador de qualidade que avalia os cursos por intermédio dos desempenhos dos estudantes a partir dos resultados do Enade.
Os valores possíveis vão de 1 ( pior situação ) a 5 ( melhor situação ).
Os cursos com apenas um ou sem nenhum concluinte participante não obtêm o Conceito Enade, ficando Sem Conceito ( SC ).
Na Tabela 5.1, são apresentadas a quantidade e a distribuição de cursos de Administração participantes do Enade/2018, por faixa de conceito e Grande Região. \textbf{Enfatiza} -se, mais uma vez, que a diferença entre os cursos tabulados, neste capítulo e no capítulo 2, corresponde aos cursos Sem Conceito, ou seja àqueles, em princípio, sem alunos concluintes que houvessem participado da prova ou àqueles que tivessem tão somente um aluno concluinte¹⁹.
Observando -se os dados da Tabela 5.1, nota -se que, dos 1.765 cursos participantes, 887 ( 50,3 \% ) classificaram -se com conceito 3, o valor modal.
Esse foi também o conceito modal em todas as regiões : Norte ( 45,8 \% ), Nordeste ( 50,3 \% ), Sudeste ( 47,4 \% ), Sul ( 57,6 \% ) e Centro-Oeste ( 50,6 \% ).
O conceito 2 foi o segundo mais frequente em nível nacional ( 21,9 \%, correspondendo a 386 cursos ), e o conceito 4, o terceiro ( 20,3 \%, correspondendo a 359 cursos ).
Houve, ainda, 79 cursos ( 4,5 \% ) que receberam conceito 5, e 28 cursos ( 1,6 \% ) que receberam conceito 1.
Dos 1.765 cursos de Administração, 26 ( 1,5 \% ) ficaram Sem Conceito ( SC ). ¹⁹ Esses cursos com somente um concluinte não foram considerados no capítulo 2.
A região Norte participou com 96 cursos ou 5,4 \% do total nacional.
Desses, 44 cursos receberam o conceito 3, o que equivale a 45,8 \% do total regional, além de ser o conceito modal na região, como já comentado.
A 33 cursos ( 34,4 \% ) atribuiu -se conceito 2, a 14 cursos ( 14,6 \% ), conceito 4, a três cursos ( 3,1 \% ), conceito 5, e a um curso, conceito 1 ( 1,0 \% ).
Um curso ficou Sem Conceito ( SC ) nessa região.
A região Nordeste participou com 324 cursos ou 18,4 \% do total nacional.
Nessa região, 163 cursos ( 50,3 \% em termos regionais ) obtiveram conceito 3, o conceito modal para a região, como já comentado.
Os conceitos 2 e 4 foram atribuídos, respectivamente, a 87 e a 55 cursos ( 26,9 \% e 17,0 \%, respectivamente ).
Ainda nessa região, 11 cursos ( 3,4 \% ) foram avaliados com o conceito 5, e apenas dois cursos ( 0,6 \% ) receberam conceito 1.
Nessa região, seis ( 1,9 \% ) cursos ficaram Sem Conceito ( SC ).
Dos 806 cursos participantes da região Sudeste ( 45,7 \% do total nacional ), o conceito 3 foi o conceito modal, como já anteriormente assinalado, tendo sido obtido por 382 cursos ( 47,4 \% ).
O conceito 1 foi atribuído a 14 cursos ( 1,7 \% ), e o conceito 2, a 163 cursos ( 20,2 \% ).
Receberam os conceitos 4 e 5, respectivamente, 186 cursos ( 23,1 \% ) e 46 cursos ( 5,7 \% ).
Nessa região, 15 ( 1,9 \% ) cursos ficaram Sem Conceito ( SC ).
Os 363 cursos da região Sul corresponderam a 20,6 \% do total nacional.
Como já indicado, houve predominância do conceito 3, atribuído a 57,6 \% dos cursos da região ( conceito modal ), ou seja, atribuído a 209 dos 363 cursos participantes na região Sul.
O conceito 2 foi atribuído a 52 cursos ( 14,3 \% ), e os conceitos 4 e 5, a 80 e 14 cursos, respectivamente, correspondendo a ( 22,0 \% e 3,9 \%, respectivamente ).
Cinco cursos receberam o conceito 1, e três ( 0,8 \% ) cursos ficaram Sem Conceito ( SC ), na região Sul.
Já dos 176 cursos participantes na região Centro-Oeste ( 10,0 \% do total nacional ), 89 ( 50,6 \% em termos regionais ) receberam conceito 3, o conceito modal, como já fora destacado.
Já seis cursos ( 3,4 \% ) foram alocados ao conceito 1, 51 cursos ( 29,0 \% ) obtiveram conceito 2, 24 ( 13,6 \% ), conceito 4 e cinco ( 2,8 \% ) receberam o conceito 5.
Além disso, apenas um ( 0,6 \% ) curso ficou Sem Conceito ( SC ) nessa região.
No Gráfico 5.1, é apresentada a distribuição cumulativa do Conceito Enade segundo Grande Região.
Em um gráfico de distribuição cumulativa, em cada valor de abscissa ( eixo x ), o valor apresentado no eixo das ordenadas ( eixo y ) é o valor acumulado da variável.
Por exemplo, as regiões Norte ( linha azul ) e Centro-Oeste ( linha laranja ) apresentam mais de 80 \% dos cursos com conceito 3 ou abaixo deste valor.
A região Nordeste ( linha verde ), por outro lado, apresenta um pouco menos de 80 \% dos cursos na mesma situação : com conceito 3 ou abaixo.
Já nas regiões Sudeste ( linha vermelha ) e Sul ( linha rosa ), a mesma situação para com conceito 3 ou abaixo é observada por menos de 74 \% dos cursos.
Em linhas gerais, em um gráfico de distribuição cumulativa, poligonais mais à esquerda representam Grandes Regiões com pior distribuição, e poligonais mais à direita, Grandes Regiões com melhores conceitos.
As regiões Norte ( linha azul ) e Centro-Oeste ( linha laranja ) apresentam os cursos com conceitos mais baixos, ao passo que as regiões Sudeste ( linha vermelha ) e Sul ( linha rosa ) apresentam as distribuições com valores maiores.
Cursos em IES situadas na região Nordeste ( linha verde ) apresentam uma situação intermediária. 5.2 CONCEITOS POR CATEGORIA ADMINISTRATIVA, POR MODALIDADE DE ENSINO E GRANDE REGIÃO Na Tabela 5.2, é apresentada a distribuição dos cursos de Administração participantes do Enade/2018, por Categoria Administrativa e Modalidade de Ensino, de acordo com os conceitos por eles alcançados, segundo as Grandes Regiões brasileiras.
Dos 1.765 cursos participantes, 248 ( 14,1 \% ) eram ministrados em Instituições Públicas, e 1.517 ( 85,9 \% ), em Instituições Privadas.
Quanto à Modalidade de Ensino, 69 cursos ( 3,9 \% ) eram a Distância, e 1.696 ( 96,1 \% ), Presenciais.
De acordo com as informações da Tabela 5.2, em termos nacionais, entre Instituições Públicas e Privadas, observa -se que, dos 79 cursos avaliados com conceito 5, 45 eram oferecidos em IES Públicas, e 34, em IES Privadas.
Dos 248 cursos participantes de IES Públicas, o conceito 4 foi o valor modal, atribuído a 95 cursos ( 38,3 \% ).
Dos demais cursos de Instituições Públicas participantes, 74 ( 29,8 \% ) cursos foram avaliados com conceito 3, 45 cursos ( 18,1 \% ) foram avaliados com conceito 5, como já comentado, 30 cursos ( 12,1 \% ), com conceito 2, e um curso ( 0,4 \% ), com conceito 1.
Três cursos ( 1,2 \% ) ficaram Sem Conceito ( SC ) nessa categoria.
Na rede privada, o conceito modal foi 3, atribuído a 813 cursos, o correspondente a 53,6 \% dos 1.517 cursos da categoria.
Dos demais cursos participantes, 27 ( 1,8 \% ) receberam conceito 1, e 356 ( 23,5 \% ), conceito 2.
O conceito 4 foi atribuído a 264 cursos ( 17,4 \% ), e o conceito 5, a 34 cursos ( 2,2 \% ).
Nessa Categoria Administrativa, 23 cursos ( 1,5 \% ) ficaram Sem Conceito ( SC ).
Dos 1.696 cursos participantes da modalidade presencial, o conceito 3 foi o valor modal, atribuído a 850 cursos ( 50,1 \% ).
Dos demais cursos nessa Modalidade de Ensino, 28 cursos ( 1,7 \% ) receberam conceito 1, 373 cursos ( 22,0 \% ) receberam conceito 2, 346 cursos ( 20,4 \% ), conceito 4, e 74 cursos ( 4,4 \% ), conceito 5.
Nessa Modalidade de Ensino, 25 cursos ( 1,5 \% ) ficaram Sem Conceito ( SC ).
Já 69 cursos foram oferecidos na modalidade a Distância.
Para esses, a moda foi o conceito 3, com 37 cursos ( 53,6 \% ).
Nessa modalidade, 13 ( 18,8 \% ) cursos foram alocados ao conceito 2, outros 13 receberam o conceito 4, e cinco obtiveram conceito 5, e somente um curso ficou Sem Conceito ( SC ).
Nenhum curso recebeu conceito 1 nessa Modalidade de Ensino.
Na análise por região, observa -se que, na região Norte, as Instituições Privadas participaram com 75 cursos ( 78,1 \% do total regional ), dos quais um ( 1,3 \% ) obteve conceito 1, um outro recebeu o conceito 5, e um ficou Sem Conceito ( SC ).
O conceito modal para as Instituições Privadas foi 3, atribuído a 36 cursos, correspondendo a 48,0 \% dessa categoria na região.
Nessa categoria e região, 28 cursos ( 37,3 \% ) foram avaliados com conceito 2, e oito cursos ( 10,7 \% ), com conceito 4.
As Instituições Públicas participaram com 21 cursos na região Norte ( 21,9 \% do total regional ), aos quais foram atribuídos conceito 2 ( cinco cursos ), conceito 3 ( oito cursos, valor modal ), conceito 4 ( seis cursos ) e conceito 5 ( dois cursos ).
Nenhum curso ficou Sem Conceito ( SC ) ou com conceito 1 nessa combinação de categoria e região.
Na região Norte, a maioria dos cursos foi oferecido na modalidade Presencial : 95 cursos ( 99,0 \% ).
Para esses, a moda foi o conceito 3, com 43 cursos.
Dos demais cursos presenciais da região, um foi alocado ao conceito 1, 33 receberam o conceito 2, 14 obtiveram conceito 4, e apenas três cursos foram alocados ao conceito 5.
Um curso ficou Sem Conceito ( SC ), nessa região.
Na região Norte, o único curso oferecido na modalidade Educação a Distância obteve o conceito 3.
Na região Nordeste, a rede privada concentrou 255 dos 324 cursos participantes, o equivalente a 78,7 \% do total da região.
Dentre os cursos oferecidos por IES Privadas no Nordeste, 142 cursos foram avaliados com conceito 3, o conceito modal.
Dois cursos receberam conceito 1, e outros dois cursos receberam conceito 5, 80 cursos obtiveram conceito 2, 23 cursos receberam conceito 4.
Nessa combinação de categoria e região, seis cursos ficaram Sem Conceito ( SC ).
As Instituições Públicas da região Nordeste participaram com 69 cursos ( 21,3 \% do total da região ).
Desses, 32 obtiveram conceito 4, o valor modal, seguidos por 21 cursos, que receberam conceito 3, além de sete cursos que obtiveram conceito 2, e nove cursos foram alocados ao conceito 5.
Nenhum curso ficou Sem Conceito ( SC ) ou recebeu o conceito 1 nessa região.
A maioria dos cursos oferecidos na região Nordeste foi na modalidade Presencial : 317 cursos ( 97,8 \% ).
Para esses, a moda foi o conceito 3, com 157 cursos.
Dos demais cursos presenciais da região, dois foram alocados ao conceito 1, 87 receberam o conceito 2, 54 obtiveram conceito 4 e 11 receberam o conceito 5.
Seis cursos ficaram Sem Conceito ( SC ) nessa região.
Na região Nordeste, apenas sete cursos foram oferecidos na modalidade Educação a Distância, seis dos quais obtiveram o conceito 3 ( conceito modal ) ; e apenas um curso que recebeu o conceito 4.
Na região Sudeste, a proporção de cursos da rede privada, 90,9 \%, foi mais \textbf{elevada} do que a registrada nas demais regiões brasileiras, correspondendo a 733 dos 806 cursos participantes.
Nessa categoria e região, o conceito modal foi 3, atribuído a 367 cursos e, 14 cursos ficaram Sem Conceito ( SC ).
Os demais foram avaliados com conceito 1 ( 13 cursos ), conceito 2 ( 156 cursos ), conceito 4 ( 160 cursos ) e conceito 5 ( 23 cursos ).
Dos 73 cursos oferecidos por Instituições Públicas na região Sudeste, o conceito modal foi 4, atribuído a 26 cursos.
Nessa região, um curso de IES Públicas ficou Sem Conceito ( SC ), enquanto um outro curso recebeu conceito 1, sete cursos obtiveram conceito 2, 15 cursos foram alocados ao conceito 3, 26 cursos receberam o conceito 4, o conceito modal, e 23 cursos ficaram com o conceito 5.
Também na região Sudeste, a maioria dos cursos foi oferecido na modalidade Presencial : 769 cursos.
Para esses, a moda foi o conceito 3, com 364 cursos.
Os demais cursos presenciais da região foram alocados aos conceitos 1 ( 14 cursos ), 2 ( 156 cursos ), 4 ( 179 cursos ) e 5 ( 42 cursos ).
Ficaram Sem Conceito ( SC ) 14 cursos.
Na região Sudeste, foram oferecidos 37 cursos na modalidade Educação a Distância que obtiveram os conceitos 2 ( sete cursos ), 3 ( 18 cursos, valor modal ), 4 ( sete cursos ) e 5 ( quatro cursos ).
Somente um curso ficou Sem Conceito ( SC ), e nenhum curso recebeu o conceito 1, nessa combinação de modalidade e região.
As Instituições Privadas concentraram 311 dos 363 cursos participantes da região Sul, 85,7 \% do total regional.
Desses, 192 cursos obtiveram conceito 3, o conceito modal.
Nesta combinação de Categoria Administrativa e Grande Região, os demais cursos receberam conceito 1 ( cinco cursos ), enquanto outros cinco obtiveram conceito 5.
Os conceitos 2 e 4 receberam 49 e 59 cursos, respectivamente.
Somente um curso ficou Sem Conceito ( SC ).
As Instituições Públicas na região Sul participaram com 52 cursos ( 14,3 \% dos cursos da região ), aos quais foram atribuídos os conceitos 2 ( três cursos ), 3 ( 17 cursos ), 4 ( 21 cursos, valor modal ), e 5 ( nove cursos ).
Nenhum curso da região Sul oferecido por IES Públicas recebeu o conceito 1, e dois cursos ficaram Sem Conceito ( SC ).
Na região Sul, também, a maioria dos cursos foi oferecida na modalidade Presencial : 346 cursos ( 95,3 \% ).
Para esses, a moda foi o conceito 3, com 200 cursos.
Dos demais cursos presenciais da região, cinco foram alocados ao conceito 1, 48 receberam o conceito 2, 76 obtiveram conceito 4, e 14 cursos foram alocados ao conceito 5.
Apenas três cursos ficaram Sem Conceito ( SC ) nessa região.
Na região Sul, foram oferecidos 17 cursos na modalidade Educação a Distância, nove dos quais obtiveram o conceito 3, valor modal, e quatro cursos receberam o conceito 2, enquanto outros quatro ficaram com o conceito 4.
Nessa combinação de modalidade e região, nenhum curso ficou Sem Conceito ( SC ) ou recebeu os conceitos 1 ou 5.
Dos demais cursos, seis receberam conceito 1, 43 obtiveram conceito 2, 14 ficaram com o conceito 4 e três foram alocados ao conceito 5.
Somente um curso ficou Sem Conceito ( SC ).
Quanto aos 33 cursos oferecidos por Instituições Públicas na região Centro-Oeste ( 18,8 \% do total regional ), 13 cursos apresentaram conceito 3, conceito modal, oito ficaram com o conceito 2, dez receberam o conceito 4, e dois obtiveram conceito 5.
Nenhum curso ficou Sem Conceito ( SC ) ou recebeu o conceito 1 nessa combinação de categoria e região.
A maioria dos cursos oferecidos na região Centro-Oeste, também, foi na modalidade Presencial : 169 cursos ( 96,0 \% ).
Para esses, a moda foi o conceito 3, com 86 cursos.
Os demais cursos presenciais da região, seis foram alocados ao conceito 1, 49 receberam o conceito 2, 23 obtiveram conceito 4 e quatro receberam o conceito 5.
Somente um curso ficou Sem Conceito ( SC ) nessa região.
Na região Centro-Oeste, apenas sete cursos foram oferecidos na modalidade Educação a Distância, três dos quais obtiveram o conceito 3 ( conceito modal ), dois receberam o conceito 2, e apenas um curso recebeu o conceito 4, enquanto um outro recebeu o conceito 5.
No Gráfico 5.2, apresenta -se a distribuição do Conceito Enade, segundo a Categoria Administrativa da IES.
Os cursos em IES Públicas ( linha azul ) apresentam uma poligonal mais à direita do que os em IES Privadas ( linha verde ) e, consequentemente, uma distribuição de Conceitos Enade com valores maiores. 5.3 CONCEITOS POR ORGANIZAÇÃO ACADÊMICA E GRANDE REGIÃO Na Tabela 5.3, encontra -se a distribuição dos conceitos atribuídos aos cursos participantes do Enade/2018, na Área de Administração, por Organização Acadêmica, segundo as Grandes Regiões brasileiras.
Dos 1.765 cursos de Administração participantes, 515 eram oferecidos em Universidades, 292, em Centros Universitários, 931, em Faculdades e 27, em CEFET/IFET.
Esta distribuição corresponde a, respectivamente, 29,2 \%, 16,5 \%, 52,7 \% e 1,5 \% dos cursos.
De acordo com os dados apresentados, de todos os 79 cursos avaliados com conceito 5, 42 eram vinculados a Universidades.
Os cursos vinculados a esse tipo de Organização Acadêmica tiveram o conceito 3 como conceito modal, atribuído a 248 cursos ( 48,2 \% ).
Os demais cursos vinculados a Universidades receberam os conceitos 1 ( três cursos ), 2 ( 53 cursos ), 4 ( 161 cursos ) e conceito 5 ( 42 cursos, como já mencionado ).
Ficaram Sem Conceito ( SC ) oito cursos.
Dos cursos em Centros Universitários, o conceito modal também foi 3, atribuído a 174 cursos ( 59,6 \% ).
Dois cursos vinculados a esse tipo de Organização Acadêmica ficaram Sem Conceito ( SC ), e os outros cursos receberam os conceitos 1 ( três cursos ), 2 ( 40 cursos ), 4 ( 62 cursos ) e 5 ( 11 cursos ).
Dos 931 cursos mantidos por Faculdades, 459 ( 49,3 \% ) receberam conceito 3, o conceito modal.
O segundo conceito mais frequentemente atribuído aos cursos oferecidos por essa categoria de IES foi o conceito 2, recebido por 291 cursos.
Na sequência, 124 cursos obtiveram conceito 4, 22 cursos, conceito 1, e 19 cursos, conceito 5.
Dos cursos oferecidos em Faculdades, 16 ficaram Sem Conceito ( SC ).
Dos 27 cursos mantidos por CEFET/IFET, 12 ( 44,4 \% ) receberam conceito 4, o conceito modal.
O segundo conceito mais frequentemente atribuído aos cursos oferecidos por essa categoria de IES foi o conceito 5, recebido por sete cursos.
Na sequência, seis cursos receberam conceito 3, e dois cursos receberam conceito 2.
Nenhum curso ficou Sem Conceito ( SC ) ou recebeu conceito 1.
Considerando -se, separadamente, as regiões brasileiras, verifica -se que, na região Norte, as Universidades participaram com 22 dos 96 cursos avaliados, dos quais dez obtiveram conceito 3 ( conceito modal ) e, na sequência, sete cursos receberam o conceito 4, três ficaram com conceito 2 e dois foram atribuídos ao conceito 5.
Nenhum curso ficou Sem Conceito ( SC ) ou recebeu o conceito 1 nessa combinação de categoria e região.
Os Centros Universitários da região Norte foram representados por 13 cursos, oito dos quais obtiveram conceitos 3 ( conceito modal ) seguidos de três que receberam o conceito 2, e um curso foi alocado ao conceito 4, enquanto um outro curso recebeu o conceito 5.
Nenhum curso ficou Sem Conceito ( SC ) ou recebeu o conceito 1 nessa combinação de categoria e região.
As Faculdades participaram com 60 cursos na região Norte, 26 dos quais receberam os conceitos 2 e 3 ( conceitos modais ), seis cursos obtiveram conceito 4 e um curso foi alocado ao conceito 1.
Um curso ficou Sem Conceito ( SC ) e nenhum obteve o conceito 5.
Já os CEFET/IFET participaram somente com um curso na região Norte, e a esse foi atribuído o conceito 2.
Na região Nordeste, as Universidades participaram com 80 dos 324 cursos da Área de Administração oferecidos.
Apenas dois desses 80 cursos ficou Sem Conceito ( SC ).
O conceito modal 3 foi atribuído a 33 cursos e o conceito modal 4, a outros 33.
Dos cursos restantes, oito obtiveram conceito 5, e quatro ficaram com conceito 2.
Nenhum curso oferecido por Universidades no Nordeste recebeu conceito 1.
Os Centros Universitários contaram com 46 cursos participantes na região Nordeste, 32 dos quais receberam o conceito modal 3.
Para os demais cursos, foram atribuídos os conceitos 2 ( seis cursos ), 4 ( sete cursos ) e 5 ( um curso ).
Nenhum curso oferecido por Centros Universitários no Nordeste ficou Sem Conceito ( SC ) ou recebeu o conceito 1.
As Faculdades foram representadas por 194 cursos na região Nordeste, 97 dos quais receberam conceito 3, o conceito modal.
Na sequência, 77 cursos obtiveram conceito 2, 13 receberam conceito 4, dois ficaram com conceito 1 e um curso foi alocado ao conceito 5.
Na região Nordeste, quatro cursos mantidos por Faculdades ficaram Sem Conceito ( SC ).
Os CEFET/IFET participaram com quatro cursos na região Nordeste, dois dos quais receberam conceito 4, um curso ficou com conceito 3, e um outro com conceito 5.
Nenhum curso ficou Sem Conceito ( SC ), ou recebeu os conceitos 1 ou 2.
Na região Sudeste, as Universidades concentraram 221 dos 806 cursos de Administração da região.
Dos cursos oferecidos em Universidades, nessa região, o conceito modal foi 3, atribuído a 111 cursos.
Os demais cursos receberam os conceitos 1 ( um curso ), 2 ( 24 cursos ), 4 ( 62 cursos ) e 5 ( 20 cursos ).
Nessa combinação de organização e região, três cursos ficaram Sem Conceito ( SC ).
Os Centros Universitários participaram com 158 cursos na região Sudeste, 84 dos quais obtiveram o conceito modal 3.
O conceito 1 foi atribuído a dois cursos, o conceito 2, a 25 cursos, o conceito 4, a 38 cursos, e o conceito 5, a sete cursos.
Apenas dois cursos ficaram Sem Conceito ( SC ) nessa combinação de Organização Acadêmica e região.
Já as Faculdades foram representadas por 409 cursos na região Sudeste, dez dos quais ficaram Sem Conceito ( SC ).
O conceito modal 3 foi atribuído a 183 cursos.
Aos demais cursos foram atribuídos os conceitos 1 ( 11 cursos ), 2 ( 113 cursos ), 4 ( 79 cursos ) e 5 ( 13 cursos ).
Os CEFET/IFET participaram com 18 cursos, sete dos quais receberam o conceito 4 ( conceito modal ), seguidos de seis cursos que obtiveram conceito 5, quatro ficaram com conceito 3, e somente um curso foi alocado ao conceito 2.
Nenhum curso ficou Sem Conceito ( SC ) ou recebeu o conceito 1.
Dos 363 cursos da região Sul, 150 eram mantidos por Universidades, para os quais o conceito modal foi 3, atribuído a 75 cursos.
Aos demais cursos foram atribuídos os conceitos 1 ( dois cursos ), 2 ( 14 cursos ), 4 ( 47 cursos ) e 5 ( nove cursos ).
Nesse tipo de organização, três cursos da região Sul ficaram Sem Conceito ( SC ).
Dos 51 cursos participantes de Centros Universitários da região Sul, nenhum ficou Sem Conceito ( SC ) ou recebeu o conceito 1.
O conceito modal 3 foi obtido por 33 cursos, seguidos por 11 cursos que receberam conceito 4, cinco obtiveram conceito 2, e dois ficaram com conceito 5.
Quanto aos 158 cursos vinculados a Faculdades na região Sul, 100 receberam o conceito modal 3.
O conceito 2 foi atribuído a 33 cursos, e o conceito 4, a 19.
Três cursos obtiveram conceito 1, enquanto outros três receberam conceito 5.
Nenhum curso ficou Sem Conceito ( SC ).
Já os CEFET/IFET participaram com quatro cursos na região Sul, três dos quais receberam conceito 4, e um curso recebeu conceito 3.
Na região Centro-Oeste, 42 dos 176 cursos oferecidos eram mantidos por Universidades.
Para este tipo de Organização Acadêmica, nenhum curso ficou Sem Conceito ( SC ) ou recebeu o conceito 1, os conceitos 2 e 4 foram atribuídos a, respectivamente, oito e 12 cursos, o conceito 3, conceito modal, foi atribuído a 19 cursos, e três cursos obtiveram conceito 5.
Os Centros Universitários da região Centro-Oeste contaram com 24 cursos, 17 dos quais receberam conceito 3, em valor modal.
Aos demais cursos foram atribuídos os conceitos 1 ( um curso ), 2 ( um curso ) e 4 ( cinco cursos ).
Nenhum curso ficou Sem Conceito ( SC ) ou recebeu o conceito 5.
Dos 110 cursos oferecidos em Faculdades na região Centro-Oeste, 53 receberam o conceito modal 3, seguidos de 42 cursos alocados ao conceito 2, sete ficaram com conceito 4, cinco obtiveram conceito 1, e dois cursos receberam conceito 5.
Nessa região, somente um curso oferecido em Faculdades ficou Sem Conceito ( SC ).
Na região Centro-Oeste, não houve registro de nenhum curso da Organização Acadêmica CEFET/IFET.
No Gráfico 5.3, apresenta -se a distribuição do Conceito Enade, segundo a Organização Acadêmica da IES.
Os cursos em CEFET/IFET ( linha laranja ) apresentam uma poligonal mais à direita do que as demais, denotando uma melhor distribuição de Conceitos Enade.
No outro extremo, encontram -se os cursos mantidos por Faculdades ( linha vermelha ), apresentando a respectiva poligonal mais à esquerda e, consequentemente, distribuição de Conceitos Enade com menores valores.
Já os cursos em Universidades ( linha azul ) e Centros Universitários ( linha verde ) apresentam distribuição de Conceitos Enade com valores intermediários.

% matched lemmas: elevado, enfatizar
\end{textsample}
